\documentclass[a4paper]{article}
\usepackage[fontset=fandol]{ctex}
\usepackage{amsmath, amssymb}
\usepackage{graphicx}
\usepackage[margin=2.5cm]{geometry}
\usepackage[most]{tcolorbox}
\usepackage{etoolbox}
\usepackage{listings}
\usepackage{hyperref}
\usepackage{booktabs}
\usepackage{subcaption}
\usepackage{float}
\usepackage{tikz}
\IfFileExists{pgfplots.sty}{\usepackage{pgfplots}\pgfplotsset{compat=1.18}}{}
\newtcolorbox{knowledgebox}[1]{enhanced, colback=blue!5!white, colframe=blue!75!black, colbacktitle=blue!75!black, coltitle=white, fonttitle=\bfseries, title=#1, attach boxed title to top left={yshift=-2mm, xshift=2mm}, boxrule=1pt, sharp corners}
\newtcolorbox{importantbox}[1]{enhanced, colback=yellow!10!white, colframe=yellow!80!black, colbacktitle=yellow!80!black, coltitle=black, fonttitle=\bfseries, title=#1, sharp corners}
\newtcolorbox{warningbox}[1]{enhanced, colback=red!5!white, colframe=red!75!black, colbacktitle=red!75!black, coltitle=white, fonttitle=\bfseries, title=#1, sharp corners}
\lstset{language=Python,basicstyle=\ttfamily\small,keywordstyle=\color{blue},stringstyle=\color{red!60!black},commentstyle=\color{green!60!black},breaklines=true,frame=single,numbers=left,numberstyle=\tiny\color{gray},captionpos=b,extendedchars=false}
\newcommand{\notetitle}{Neural Network Language Model}
\newcommand{\noteauthors}{Codex 基于课程页、公开视频字幕、coverage ledger 与补充官方材料整理}
\newcommand{\notedate}{\today}
\newcommand{\videochannel}{WAVLab}
\newcommand{\videopublishdate}{2023-11-01}
\newcommand{\videoduration}{unknown}
\newcommand{\videourl}{https://www.youtube.com/watch?v=JmDeT8X9Izg}
\newcommand{\videocoverpath}{cover.jpg}
\begin{document}
\begin{titlepage}
\centering
{\Large 课程笔记\par}
\vspace{1.2cm}
{\huge\bfseries \notetitle\par}
\vspace{0.8cm}
{\large \noteauthors\par}
\vspace{0.3cm}
{\large \notedate\par}
\vspace{1.2cm}
\ifdefempty{\videocoverpath}{{\small 公开视频封面不可用。\par}}{\includegraphics[width=0.82\textwidth,height=0.45\textheight,keepaspectratio]{\videocoverpath}\par}
\vfill
\begin{tcolorbox}[width=0.9\textwidth, colback=black!2!white, colframe=black!60, sharp corners]
\textbf{视频作者/频道}：\videochannel\par
\textbf{发布日期}：\videopublishdate\par
\textbf{视频时长}：\videoduration\par
\textbf{视频链接}：\href{\videourl}{\nolinkurl{\videourl}}
\end{tcolorbox}
\end{titlepage}
\tableofcontents
\newpage
\section{CNNs for sequence inputs}
根据课程页主题、公开视频字幕以及本讲的 coverage ledger，CNNs for sequence inputs 这一部分主要围绕 Neural Network Language Model | Neural network language model；从前一讲的 acoustic modeling 回到序列建模，说明 CNN / spectrogram 输入如何延伸到文本或声学序列。 展开。
写作上保持 coverage-first 原则：先把老师在这一阶段真正讲过的问题设定、模型部件、训练或推理逻辑讲清，再压缩成适合复习的结构。

\begin{figure}[H]
\centering
\includegraphics[width=0.92\textwidth]{figures/coverage_map.png}
\caption{本讲 coverage map：按 coverage\_units.jsonl 归纳出的核心主题与章节映射。}
\end{figure}

\begin{figure}[H]
\centering
\includegraphics[width=0.92\textwidth]{figures/timeline.png}
\caption{本讲 timeline：根据字幕时间范围归纳的主题推进顺序。}
\end{figure}

\subsection{Neural Network Language Model}
从证据位置看，这一部分对应的课程证据区间是 2023-11-01。课程在这里强调的核心内容可以概括为：Neural Network Language Model | Neural network language model。
从课程组织上看，这一部分承担的是把局部概念嵌回整条识别流水线中的作用，避免读者把公式或模块孤立地记忆。

在本章结构中，这一单元被并入“CNNs for sequence inputs”，目的是把同一阶段的定义、机制和权衡集中讲清。

\subsection{从前一讲的 acoustic modeling 回到序列建模，说明 CNN / spectrogram 输入如何延伸到文本或声学序列}
从证据位置看，这一部分对应的课程证据区间是 00:00:00,960--00:07:03,319。课程在这里强调的核心内容可以概括为：从前一讲的 acoustic modeling 回到序列建模，说明 CNN / spectrogram 输入如何延伸到文本或声学序列。。
从课程组织上看，这一部分承担的是把局部概念嵌回整条识别流水线中的作用，避免读者把公式或模块孤立地记忆。

在本章结构中，这一单元被并入“CNNs for sequence inputs”，目的是把同一阶段的定义、机制和权衡集中讲清。

\subsection{本章小结}
CNNs for sequence inputs 这一部分把 Neural Network Language Model 中相邻的教学单元串成了一条连续叙事：既交代了这一阶段要解决的问题，也交代了方法、推理或工程权衡为什么要这样组织。

\section{CNN design choices}
根据课程页主题、公开视频字幕以及本讲的 coverage ledger，CNN design choices 这一部分主要围绕 CNN 设计、batch normalization、activation、residual / gating 等结构细节。 展开。
写作上保持 coverage-first 原则：先把老师在这一阶段真正讲过的问题设定、模型部件、训练或推理逻辑讲清，再压缩成适合复习的结构。

\subsection{CNN 设计、batch normalization、activation、residual / gating 等结构细节}
从证据位置看，这一部分对应的课程证据区间是 00:06:59,680--00:14:30,600。课程在这里强调的核心内容可以概括为：CNN 设计、batch normalization、activation、residual / gating 等结构细节。。
从课程组织上看，这一部分承担的是把局部概念嵌回整条识别流水线中的作用，避免读者把公式或模块孤立地记忆。

在本章结构中，这一单元被并入“CNN design choices”，目的是把同一阶段的定义、机制和权衡集中讲清。

\subsection{本章小结}
CNN design choices 这一部分把 Neural Network Language Model 中相邻的教学单元串成了一条连续叙事：既交代了这一阶段要解决的问题，也交代了方法、推理或工程权衡为什么要这样组织。

\section{Recurrent sequence models}
根据课程页主题、公开视频字幕以及本讲的 coverage ledger，Recurrent sequence models 这一部分主要围绕 time-delay / recurrent 视角下的序列模型，以及从局部卷积到递归建模的过渡。 展开。
写作上保持 coverage-first 原则：先把老师在这一阶段真正讲过的问题设定、模型部件、训练或推理逻辑讲清，再压缩成适合复习的结构。

\subsection{time-delay / recurrent 视角下的序列模型，以及从局部卷积到递归建模的过渡}
从证据位置看，这一部分对应的课程证据区间是 00:14:27,199--00:22:29,919。课程在这里强调的核心内容可以概括为：time-delay / recurrent 视角下的序列模型，以及从局部卷积到递归建模的过渡。。
从课程组织上看，这一部分承担的是把局部概念嵌回整条识别流水线中的作用，避免读者把公式或模块孤立地记忆。

在本章结构中，这一单元被并入“Recurrent sequence models”，目的是把同一阶段的定义、机制和权衡集中讲清。

\subsection{本章小结}
Recurrent sequence models 这一部分把 Neural Network Language Model 中相邻的教学单元串成了一条连续叙事：既交代了这一阶段要解决的问题，也交代了方法、推理或工程权衡为什么要这样组织。

\section{Bidirectional recurrence}
根据课程页主题、公开视频字幕以及本讲的 coverage ledger，Bidirectional recurrence 这一部分主要围绕 bidirectional networks、双向信息流与长时依赖的建模问题。 展开。
写作上保持 coverage-first 原则：先把老师在这一阶段真正讲过的问题设定、模型部件、训练或推理逻辑讲清，再压缩成适合复习的结构。

\subsection{bidirectional networks、双向信息流与长时依赖的建模问题}
从证据位置看，这一部分对应的课程证据区间是 00:22:26,559--00:29:55,760。课程在这里强调的核心内容可以概括为：bidirectional networks、双向信息流与长时依赖的建模问题。。
从课程组织上看，这一部分承担的是把局部概念嵌回整条识别流水线中的作用，避免读者把公式或模块孤立地记忆。

在本章结构中，这一单元被并入“Bidirectional recurrence”，目的是把同一阶段的定义、机制和权衡集中讲清。

\subsection{本章小结}
Bidirectional recurrence 这一部分把 Neural Network Language Model 中相邻的教学单元串成了一条连续叙事：既交代了这一阶段要解决的问题，也交代了方法、推理或工程权衡为什么要这样组织。

\section{LSTM and memory}
根据课程页主题、公开视频字幕以及本讲的 coverage ledger，LSTM and memory 这一部分主要围绕 LSTM / memory 机制如何缓解长程依赖，并过渡到 neural language model 的概率视角。 展开。
写作上保持 coverage-first 原则：先把老师在这一阶段真正讲过的问题设定、模型部件、训练或推理逻辑讲清，再压缩成适合复习的结构。

\subsection{LSTM / memory 机制如何缓解长程依赖，并过渡到 neural language model 的概率视角}
从证据位置看，这一部分对应的课程证据区间是 00:29:53,240--00:37:17,960。课程在这里强调的核心内容可以概括为：LSTM / memory 机制如何缓解长程依赖，并过渡到 neural language model 的概率视角。。
从课程组织上看，这一部分承担的是把局部概念嵌回整条识别流水线中的作用，避免读者把公式或模块孤立地记忆。

在本章结构中，这一单元被并入“LSTM and memory”，目的是把同一阶段的定义、机制和权衡集中讲清。

\subsection{本章小结}
LSTM and memory 这一部分把 Neural Network Language Model 中相邻的教学单元串成了一条连续叙事：既交代了这一阶段要解决的问题，也交代了方法、推理或工程权衡为什么要这样组织。

\section{Neural language model formulation}
根据课程页主题、公开视频字幕以及本讲的 coverage ledger，Neural language model formulation 这一部分主要围绕 neural language model 的表示、embedding / inner product 以及序列概率建模，并与 n-gram / 计数式 LM 做方法对比。 展开。
写作上保持 coverage-first 原则：先把老师在这一阶段真正讲过的问题设定、模型部件、训练或推理逻辑讲清，再压缩成适合复习的结构。

\subsection{neural language model 的表示、embedding / inner product 以及序列概率建模，并与 n-gram / 计数式 LM 做方法对比}
从证据位置看，这一部分对应的课程证据区间是 00:37:13,839--00:44:42,559。课程在这里强调的核心内容可以概括为：neural language model 的表示、embedding / inner product 以及序列概率建模，并与 n-gram / 计数式 LM 做方法对比。。
课程把 n-gram language model 放在“可解释但近似”的框架里讨论，重点是条件独立假设、估计、平滑以及搜索中的作用。

在本章结构中，这一单元被并入“Neural language model formulation”，目的是把同一阶段的定义、机制和权衡集中讲清。
\[
P(w_1,\dots,w_T) = \prod_{t=1}^{T} P(w_t \mid w_1,\dots,w_{t-1}) \approx \prod_{t=1}^{T} P(w_t \mid w_{t-n+1},\dots,w_{t-1})
\]

符号说明如下：
\begin{itemize}
\item $w_t$：第 $t$ 个词；
\item $T$：句子长度；
\item $n$：$n$-gram 的上下文窗口长度。
\end{itemize}

\subsection{本章小结}
Neural language model formulation 这一部分把 Neural Network Language Model 中相邻的教学单元串成了一条连续叙事：既交代了这一阶段要解决的问题，也交代了方法、推理或工程权衡为什么要这样组织。

\section{Transformer masking and implementation}
根据课程页主题、公开视频字幕以及本讲的 coverage ledger，Transformer masking and implementation 这一部分主要围绕 Transformer / masking / implementation pitfalls，作为神经语言模型的现代收束。 展开。
写作上保持 coverage-first 原则：先把老师在这一阶段真正讲过的问题设定、模型部件、训练或推理逻辑讲清，再压缩成适合复习的结构。

\subsection{Transformer / masking / implementation pitfalls，作为神经语言模型的现代收束}
从证据位置看，这一部分对应的课程证据区间是 00:44:40,200--00:51:55,240。课程在这里强调的核心内容可以概括为：Transformer / masking / implementation pitfalls，作为神经语言模型的现代收束。。
从课程组织上看，这一部分承担的是把局部概念嵌回整条识别流水线中的作用，避免读者把公式或模块孤立地记忆。

在本章结构中，这一单元被并入“Transformer masking and implementation”，目的是把同一阶段的定义、机制和权衡集中讲清。

\subsection{本章小结}
Transformer masking and implementation 这一部分把 Neural Network Language Model 中相邻的教学单元串成了一条连续叙事：既交代了这一阶段要解决的问题，也交代了方法、推理或工程权衡为什么要这样组织。

\section{Implementation checks and recap}
根据课程页主题、公开视频字幕以及本讲的 coverage ledger，Implementation checks and recap 这一部分主要围绕 实现检查、masking 细节与整讲收尾。 展开。
写作上保持 coverage-first 原则：先把老师在这一阶段真正讲过的问题设定、模型部件、训练或推理逻辑讲清，再压缩成适合复习的结构。

\subsection{实现检查、masking 细节与整讲收尾}
从证据位置看，这一部分对应的课程证据区间是 00:51:53,160--00:59:08,319。课程在这里强调的核心内容可以概括为：实现检查、masking 细节与整讲收尾。。
从课程组织上看，这一部分承担的是把局部概念嵌回整条识别流水线中的作用，避免读者把公式或模块孤立地记忆。

在本章结构中，这一单元被并入“Implementation checks and recap”，目的是把同一阶段的定义、机制和权衡集中讲清。

\subsection{本章小结}
Implementation checks and recap 这一部分把 Neural Network Language Model 中相邻的教学单元串成了一条连续叙事：既交代了这一阶段要解决的问题，也交代了方法、推理或工程权衡为什么要这样组织。

\section{总结与延伸}
Neural Network Language Model 这一讲在整门《Speech Recognition and Understanding》课程中的位置，是把 Neural Network Language Model 讲成一个可复用的建模模块。
从教材化角度看，本讲最重要的不是记住单一术语，而是理解它如何嵌入完整的 automatic speech recognition pipeline：输入语音如何被表示、模型如何被训练、解码如何得到最终转写，以及这些设计分别带来什么假设和权衡。
由于公开源以课程页和公开视频字幕为主、缺少官方 slide PDF，本章在正文里尽量把术语、方法和工程语境解释完整，并把缺失项留在 omission 和 manifest 里，而不是凭空补写。
\end{document}
